\documentclass[10pt,a4paper]{amsart}
\usepackage[margin=19mm]{geometry}
\usepackage{amsmath,amssymb,mathtools}
\usepackage[scr=rsfs,cal=euler]{mathalfa}
\usepackage{xfrac}
\usepackage[unicode,hidelinks,bookmarks=false]{hyperref}
\newcommand{\ie}{i.e.\ }
\newcommand{\cf}{cf.\ }
\newcommand{\resp}{respectively\ }
\newcommand{\Subsection}{\subsection}
\providecommand{\nobreakdash}{\nobreak\hbox{-}}
\setlength{\emergencystretch}{2em}
\multlinegap0pt
\usepackage{tikz}
\usepackage{amssymb}
\newtheorem{lemma}{Lemma}
\newtheorem{theorem}{Theorem}
\newcommand{\C}{{\mathbb C}}
\title{Reducibility of linear representations, free ideals, and Kippenhahn's conjecture}
\author{Michael Stessin, Rongwei Yang}
\date{Source: arXiv:2608.14194v1 (CC BY 4.0)}
\begin{document}
\maketitle

\noindent\emph{Source and licence.} Reducibility of linear representations, free ideals, and Kippenhahn's conjecture. Version arXiv:2608.14194v1 (CC BY 4.0), distributed by arXiv under the Creative Commons Attribution 4.0 International licence, http://creativecommons.org/licenses/by/4.0/, as recorded in the arXiv licence metadata for this exact version and shown on the abstract page https://arxiv.org/abs/2608.14194v1. Full licence notice: \texttt{licenses/arxiv-2608.14194-CC-BY-4.0.txt}.\par\smallskip

\noindent\textit{Editorial scope.} Counted unit: the article's theorem lcyclic (source0.tex lines 344--346). The article's complete original proof of it is delivered with it (source0.tex lines 347--372, one \texttt{proof} environment of 26 lines). No mathematics is written, repaired or re-derived: the statement and the proof are the article's own lines, in the article's own notation, and no retained line is substituted or re-flowed.  Retained uncounted as the source-local context the proof needs: lines 338--340.  Nothing was omitted from any retained span, so the package carries no omission record.  No retained line contains a citation, so the wrapper carries no bibliography and no bibliography entry.  No retained line contains a figure environment, an image-inclusion command or drawing code, so no image is delivered and the package declares no asset dependency.  Editorial formatting only, in addition: an AMS \texttt{amsart} preamble that restates the article's own declarations and macros used by the retained lines, namely the environments \texttt{lemma}, \texttt{theorem} on separate counters, together with the standard packages those lines need.  The retained environments print in this document as Lemma 1, Theorem 1 in order of appearance, whereas the article prints the same items with the numbering of its own full document.  The complete native source of the article as distributed in the arXiv e-print for this exact version is bundled unchanged as \texttt{sources/207644/source0.tex}.
\begin{lemma}\label{minip}
Suppose $q(z_0,z)$ is a degree-$d$ minimal polynomial of tuple $A$. Then for each fixed $w\in \C^n, w\neq 0$, the minimal polynomial of $A_*(w)$ has degree no more than $d$.
\end{lemma}

\begin{theorem}\label{lcyclic}
The characteristic polynomial $Q_A$ is minimal if and only if $A$ is linearly cyclic.
\end{theorem}

\begin{proof}
If $A$ is linearly cyclic, then there exists a vector $w\in \C^n$ for which $A_*(w)$ is cyclic, implying that the minimal polynomial for $A_*(w)$ coincides with its characteristic polynomial which has degree $k$. It follows from Lemma \ref{minip} that the minimal polynomial of $A$ has degree $k$, showing that $Q_A$ is minimal.

To prove the necessity, we show that if tuple $A$ is not linearly cyclic then $Q_A$ is not minimal.
First, we let $d(z)$ denote the degree of the minimal polynomial for $A_*(z)$ and set \[r_A:=\max \{d(z): z\in \C^n\}.\] For convenience, we shall write $r_A$ often as $r$. If $d(w)=k$ for some $w\in \C^n$, then $A_*(w)$ is cyclic, contradicting the assumption that $A$ is not linearly cyclic. Thus, we assume $r=d(w)<k$. This implies that the set $\{I, A_*(w), ..., A^{r-1}_*(w)\}$ is linearly independent. For any matrix $T\in M_k(\C)$, we let vec$(T)$ denote the vector in $\C^{k^2}$ in which the $j$th block is the $j$th column of $T$. Then the $k^2\times r$ matrix 
\[M(w):=\begin{pmatrix} \text{vec}(I) & \text{vec}(A_*(w)) & \cdots & \text{vec}(A^{r-1}_*(w))\end{pmatrix}\] has rank $r$, and consequently the matrix $M^T(w)M(w)$ is a $r\times r$ invertible matrix. It follows that there is an open neighborhood of $w$ on which $M^T(z)M(z)$ is invertible. Therefore, the set $\{z\in \C^n: d(z)=r\}$ is open. In fact, more information can be derived. Suppose 
\[p(z_0,w)=z_0^r-a_{r-1}(w)z_0^{r-1}-\cdots -a_0(w)=\prod_{j=1}^r(z_0-\lambda_j(w))\] is the minimal polynomial for $A_*(w)$. Then each $\lambda_j(w)$ is an eigenvalue of $A_*(w)$. Since $A_*(tw)=tA_*(w), t\in \C$, we have $\lambda_j(tw)=t\lambda_j(w)$. Thus the functions $a_i(w)$ are homogeneous of degree $r-i$. We claim that $a_i$ is a polynomial for every $0\leq i\leq r-1$. 

Indeed, since $p(A_*(w),w)=0$, it follows that
\[M(w)\begin{pmatrix} a_0(w)\\ a_1(w) \\ \vdots \\a_{r-1}(w)\end{pmatrix}=\text{vec}(A^{r}_*(w)),\] and consequently
\begin{equation}\label{coeff}
\begin{pmatrix} a_0(w)\\ a_1(w) \\ \vdots \\a_{r-1}(w)\end{pmatrix}=\left(M^T(w)M(w)\right)^{-1}M^T(w)\text{vec}(A^r_*(w)).
\end{equation}
Equation \eqref{coeff} leads to the following two facts:

\vspace{1mm}

1) The set \[V:=\{z\in \C^n: d(z)\leq r-1\}=\{z\in \C^n: \det \left(M^T(z)M(z)\right)=0\},\] which is an algebraic variety in $\C^n$ of dimension $n-1$.

2) The coefficients $a_0(w), ..., a_{r-1}(w)$ are rational functions in $\C^n$ with possible singularities in $V$.

\vspace{1mm}

Furthermore, since each $a_j(w)$ is a symmetric polynomial in the eigenvalues $\lambda_1(w), ..., \lambda_r(w)$ of $A_*(w)$, it is bounded in a neighborhood of every point in $V$. Riemann extension theorem then implies that $a_j$ is in fact a polynomial. This shows that $p(z_0,z)$ is a minimal polynomial for tuple $A$. Since $r<k$, the characteristic polynomial $Q_A$ is thus not a minimal polynomial, completing the proof of the theorem.

\end{proof}

\end{document}
